\documentclass[11pt,a4paper]{article}

\usepackage{fontspec}
\usepackage{polyglossia}
\setdefaultlanguage{russian}
\setotherlanguage{english}
\setmainfont{CMU Serif}
\setsansfont{CMU Sans Serif}
\setmonofont{CMU Typewriter Text}

\usepackage{geometry}
\geometry{top=2cm, bottom=2.5cm, left=2.5cm, right=2.5cm}

\usepackage{amsmath,amssymb,amsthm}
\usepackage{booktabs}
\usepackage{longtable}
\usepackage{float}
\usepackage{xcolor}
\usepackage{hyperref}
\hypersetup{colorlinks, urlcolor=blue, citecolor=blue, linkcolor=blue}
\usepackage{titlesec}
\usepackage{enumitem}
\usepackage{caption}
\captionsetup{font=small, labelfont=bf}

\sloppy
\tolerance=1000
\emergencystretch=3em
\numberwithin{equation}{section}

\newtheorem{theorem}{Теорема}[section]
\newtheorem{proposition}[theorem]{Утверждение}
\newtheorem{remark}[theorem]{Замечание}

\titleformat{\section}{\Large\bfseries}{\thesection}{0.8em}{}
\titleformat{\subsection}{\large\bfseries}{\thesubsection}{0.7em}{}

\begin{document}

\begin{center}
    {\LARGE\bf Модель MEML 75}\\[0.2cm]
    {\large Феноменологическая параметризация фермионного сектора\\
    Стандартной модели на орбифолде $T^2/\mathbb{Z}_3$}\\[0.6cm]
    {\bf Юров Д. В.}\\[0.15cm]
{\small Независимый исследователь}\\[0.1cm]
{\small \texttt{Yurov@yandex.ru}}\\[0.15cm]
{\small ORCID: \texttt{0009-0006-5467-7510}}\\[0.3cm]
    {\it Финальная ревизия v3, 2026}
\end{center}

\vspace{0.4cm}

\begin{abstract}
\noindent
Мы представляем финальную ревизию модели MEML 75 ---
феноменологической параметризации фермионного сектора
Стандартной модели (СМ) на шестимерном орбифолде
$T^2/\mathbb{Z}_3$. Модель описывает шесть отношений
масс фермионов, матрицу CKM и матрицу PMNS, включая
CP-нарушение, через геометрию компактификации.

Проведено одиннадцать независимых тестов. Ранг
якобиана наблюдаемых по параметрам равен 4 в PMNS
и 4 в CKM; это совпадает с числом наблюдаемых в каждой
матрице. Число физически независимых параметров в
фермионном секторе --- 14 (6 масс + 4 CKM + 4 PMNS).
Устойчиво воспроизводится только иерархия масс
нейтрино (NO); остальные нейтринные наблюдаемые
($\Delta m^2_{21}/\Delta m^2_{31}$, $\delta_{CP}$,
$\Sigma m_\nu$, $m_{\beta\beta}$) оказываются
недоопределёнными и не могут считаться
предсказаниями. Полный перебор дискретных
$\mathbb{Z}_3$-фаз (531\,441 комбинация) не
воспроизводит PMNS с точностью лучше $10^{-3}$,
что указывает на необходимость непрерывных фаз.
Модель \textbf{не является теорией всего}:
гравитация не квантована, калибровочный и
хиггсовский секторы не выведены.
\end{abstract}

\vspace{0.3cm}
\tableofcontents
\vspace{0.5cm}

% ============================================================
\section{Введение}
% ============================================================

Стандартная модель физики элементарных частиц
содержит 19 свободных параметров, из которых 13
относятся к фермионному сектору: девять масс,
три угла смешивания кварков и фаза CP, а также
аналогичные параметры в лептонном секторе после
открытия нейтринных осцилляций.

Происхождение иерархии масс и структуры матриц
смешивания остаётся одной из центральных
нерешённых задач. Один из подходов состоит в том,
что иерархия возникает из геометрии дополнительных
измерений~\cite{ArkaniHamed1999}. В моделях с
локализованными волновыми функциями
(``split fermions'') массы генерируются перекрытием
волновых функций фермионов и хиггсовского поля
в дополнительных измерениях.

Модель MEML 75, представленная в настоящей работе,
следует этой философии. Она использует шестимерный
орбифолд $T^2/\mathbb{Z}_3$ в качестве
геометрической базы и параметризует фермионный
сектор СМ через локализацию волновых функций
в трёх неподвижных точках орбифолда.

\subsection{Цели и ограничения}

Настоящая работа \emph{не} претендует на статус
теории всего. Её цели:
\begin{enumerate}[leftmargin=1.5em,itemsep=0.1em]
\item Показать, что геометрия $T^2/\mathbb{Z}_3$
      позволяет параметризовать фермионный сектор СМ.
\item Честно разделить строго вычисляемые и
      подгоночные величины.
\item Установить структурные пределы модели через
      независимые тесты.
\item Сформулировать проверяемые предсказания
      и дорожную карту.
\end{enumerate}

Модель \emph{не} включает квантовую гравитацию,
калибровочный сектор СМ, хиггсовский сектор,
инфляцию, бариогенезис, тёмную материю.

% ============================================================
\section{Геометрия орбифолда}
% ============================================================

\subsection{Тор и его модуль}

Рассмотрим двумерный тор $T^2 = \mathbb{C}/\Lambda$,
порождённый векторами $1$ и $\tau$ с комплексным
модулем:
\begin{equation}
\tau = e^{i\pi/3} = \frac{1}{2} + i\frac{\sqrt{3}}{2}.
\end{equation}
Координата $z = x + \tau y$, $x,y \in \mathbb{R}$,
метрика:
\begin{equation}
ds^2 = dx^2 + dx\,dy + dy^2.
\end{equation}

\subsection{Орбифолд $T^2/\mathbb{Z}_3$}

Группа $\mathbb{Z}_3$ действует как
$z \to \omega z$, $\omega = e^{2\pi i/3}$.
Неподвижные точки:
\begin{equation}
z_1 = (0,0), \quad z_2 = \left(\tfrac{1}{3},\tfrac{1}{3}\right),
\quad z_3 = \left(\tfrac{2}{3},\tfrac{2}{3}\right).
\label{eq:fixed_points}
\end{equation}
Ровно три неподвижные точки --- геометрический
источник трёх поколений фермионов.

\subsection{Расстояние}

Евклидово расстояние на торе:
\begin{equation}
d(z_1,z_2) = \min_{m,n\in\mathbb{Z}}
\sqrt{(x_1-x_2+m)^2 + (x_1-x_2+m)(y_1-y_2+n) + (y_1-y_2+n)^2}.
\end{equation}
В расчётах ограничиваемся $m,n \in [-2,2]$.

% ============================================================
\section{Космологический блок}
% ============================================================

\subsection{Входные параметры}

Пять входных величин:
\begin{align}
\alpha_G &= \frac{2}{3\pi} = 0.21220659, \\
\varepsilon &= 2.860924 \times 10^{-4}, \\
\Lambda_{\text{QCD}} &= 0.075~\text{ГэВ}, \\
M_{Pl}^{\text{red}} &= 2.435 \times 10^{18}~\text{ГэВ}, \\
I_{\text{target}} &= 0.3653.
\end{align}

\subsection{Вычисляемые величины}

\begin{align}
m_{KK} &= \frac{M_{Pl}^{\text{red}}}{\sqrt{I_{\text{target}}}
\cdot \exp(2\pi/3\alpha_G)} = 2.0838 \times 10^{14}~\text{ГэВ}, \\
M_6 &= (3(M_{Pl}^{\text{red}})^2 m_{KK}^2)^{1/4}
= 2.9646 \times 10^{16}~\text{ГэВ}, \\
f_a &= M_{Pl}^{\text{red}}\varepsilon^2
= 1.9930 \times 10^{11}~\text{ГэВ}, \\
m_a &= \Lambda_{\text{QCD}}^2/f_a
= 28.22~\text{мкЭВ}, \\
\nu_a &= m_a c^2/h = 6.824~\text{ГГц}.
\end{align}

\begin{remark}
Входные параметры ($\alpha_G$, $\varepsilon$,
$\Lambda_{\text{QCD}}$, $M_{Pl}^{\text{red}}$,
$I_{\text{target}}$) не выведены из геометрии.
Утверждения ранних ревизий о ``выводе из первых
принципов'' необоснованы.
\end{remark}

% ============================================================
\section{Фермионный сектор}
% ============================================================

\subsection{Волновые функции}

Фермион $f$ локализован вблизи точки $z_i$ с гауссовой
шириной $\sigma_f^{(i)}$:
\begin{equation}
\Psi_f^{(i)}(z) \propto \exp\left(-\frac{|z-z_i|^2}{2(\sigma_f^{(i)})^2}\right).
\end{equation}
Хиггс локализован в $z_H$ с $\sigma_H = 0.15$.

\subsection{Механизм Фроггатта--Нильсена}

Для сокращения числа параметров:
\begin{equation}
\sigma_f^{(i)} = \exp(\log b_f + q_i d_f), \quad q_i = (2,1,0).
\end{equation}

\subsection{Матрица масс}

\begin{equation}
M_{ij} = y \int_{T^2/\mathbb{Z}_3}
\Psi_L^{(i)*}(z)\,\Psi_H(z)\,\Psi_R^{(j)}(z)\,d^2z
\cdot e^{i\phi_{ij}}.
\end{equation}
Физические массы --- собственные значения
$M M^\dagger$. Для CKM: $V = U_u^\dagger U_d$.
Для PMNS используется seesaw типа~I:
\begin{equation}
m_\nu = m_D M_R^{-1} m_D^T,
\end{equation}
с симметричной $M_R$:
$(M_R)_{ij} = \sqrt{M_i M_j}\,e^{i\phi_{ij}}$.

% ============================================================
\section{Численные результаты}
% ============================================================

\subsection{Иерархия масс}

\begin{table}[H]
\centering
\caption{Отношения масс. Сравнение с PDG 2024.}
\label{tab:masses}
\begin{tabular}{lccc}
\toprule
Сектор & Отношение & Рассчитано & Наблюдаемое \\
\midrule
Лептоны  & $m_\mu/m_\tau$ & $5.9464 \times 10^{-2}$ & $5.9464 \times 10^{-2}$ \\
         & $m_e/m_\tau$   & $2.8759 \times 10^{-4}$ & $2.8759 \times 10^{-4}$ \\
\midrule
$u$-кварки & $m_c/m_t$   & $7.3512 \times 10^{-3}$ & $7.3512 \times 10^{-3}$ \\
           & $m_u/m_t$   & $1.2502 \times 10^{-5}$ & $1.2503 \times 10^{-5}$ \\
\midrule
$d$-кварки & $m_s/m_b$   & $2.2344 \times 10^{-2}$ & $2.2344 \times 10^{-2}$ \\
           & $m_d/m_b$   & $1.1172 \times 10^{-3}$ & $1.1172 \times 10^{-3}$ \\
\bottomrule
\end{tabular}
\end{table}

Loss: $\mathcal{L}_{\text{lept}} \sim 10^{-21}$,
$\mathcal{L}_u \sim 10^{-19}$, $\mathcal{L}_d \sim 10^{-21}$.

\subsection{CKM}

\begin{table}[H]
\centering
\caption{Параметры CKM.}
\begin{tabular}{lccc}
\toprule
Параметр & Рассчитано & PDG 2024 & Ошибка \\
\midrule
$\theta_{12}$ & $13.003^\circ$ & $13.04^\circ$ & $0.3\%$ \\
$\theta_{13}$ & $0.211^\circ$  & $0.201^\circ$ & $5.2\%$ \\
$\theta_{23}$ & $2.397^\circ$  & $2.38^\circ$  & $0.7\%$ \\
$J$ & $3.076\times 10^{-5}$ & $3.08\times 10^{-5}$ & $0.1\%$ \\
\bottomrule
\end{tabular}
\end{table}

Loss: $\mathcal{L}_{\text{CKM}} = 3.84 \times 10^{-9}$.

\subsection{PMNS}

\begin{table}[H]
\centering
\caption{Параметры PMNS (NO, $\delta_{CP} = 197^\circ$).}
\begin{tabular}{lccc}
\toprule
Параметр & Рассчитано & PDG 2024 & Ошибка \\
\midrule
$\theta_{12}$ & $33.01^\circ$ & $33.41^\circ$ & $1.2\%$ \\
$\theta_{13}$ & $8.57^\circ$  & $8.54^\circ$  & $0.3\%$ \\
$\theta_{23}$ & $48.36^\circ$ & $49.10^\circ$ & $1.5\%$ \\
$|J|$         & $0.0096$ & $0.0097$ & $0.9\%$ \\
\bottomrule
\end{tabular}
\end{table}

Loss: $\mathcal{L}_{\text{PMNS}} = 7.66 \times 10^{-7}$.

% ============================================================
\section{Число физических параметров}
% ============================================================

Проведено три независимых теста для определения
числа физически независимых параметров.

\subsection{SVD якобиана наблюдаемых}

Якобиан $J_{ij} = \partial O_i/\partial p_j$,
$O_i$ --- 9 элементов $|U_{\text{PMNS}}|$ или
$|V_{\text{CKM}}|$. Сингулярные значения:
\begin{itemize}[leftmargin=1.5em,itemsep=0.05em]
\item PMNS: $\sigma_k = (10.5, 1.46, 1.02, 0.41,
      7.4\times 10^{-9},\ldots)$. \textbf{Ранг = 4.}
\item CKM: $\sigma_k = (0.964, 0.667, 0.191, 0.172,
      8.3\times 10^{-10},\ldots)$. \textbf{Ранг = 4.}
\end{itemize}
Разрыв $\sigma_3/\sigma_4 \sim 10^7$--$10^8$.

\subsection{Разрыв спектра Гессиана}

Численное вычисление $\mathcal{H}_{ij} = \partial^2 \mathcal{L}/\partial p_i \partial p_j$
в минимуме:
\begin{itemize}[leftmargin=1.5em,itemsep=0.05em]
\item PMNS: разрыв на позиции 4, отношение $|\lambda|$ --- $1.2\times 10^4$.
\item CKM: разрыв на позиции 4, отношение $6.7\times 10^4$.
\end{itemize}

\subsection{Тест сдвига}

Сдвиг каждого параметра на $\delta = 0.3$ с
последующей реоптимизацией. Даёт заниженные
числа (PMNS: 2 из 24, CKM: 0 из 18), поскольку
физические направления --- линейные комбинации
координат.

\subsection{Вывод}

Два согласованных метода (SVD и Гессиан) дают
\textbf{4 физических направления в PMNS и 4 в CKM},
что совпадает с числом наблюдаемых. Число
физически независимых параметров в фермионном
секторе:
\begin{equation}
N_{\text{phys}} = 6_{\text{массы}} + 4_{\text{CKM}} + 4_{\text{PMNS}} = 14.
\end{equation}
Оставшиеся $\sim 40$ параметров калибровочны.

% ============================================================
\section{Нейтринный сектор: предсказания и ограничения}
% ============================================================

Loss-функция фита PMNS использует только углы.
Иерархия, $\delta_{CP}$, $\Sigma m_\nu$, $m_{\beta\beta}$
и $\Delta m^2_{21}/\Delta m^2_{31}$ не участвуют
в фите. Проведён тест их устойчивости: 30 независимых
прогонов фита.

\subsection{Устойчивые величины}

\begin{itemize}[leftmargin=1.5em,itemsep=0.1em]
\item \textbf{Иерархия: NO}. Воспроизводится во всех
      прогонах независимо от seed.
\item $\Sigma m_\nu \approx 0.05$--$0.08$~эВ
      ($< 0.12$~эВ, Planck 2018). Согласуется с
      космологическими ограничениями.
\end{itemize}

\subsection{Неустойчивые величины}

\begin{itemize}[leftmargin=1.5em,itemsep=0.1em]
\item $\delta_{CP}$: разные минимумы дают $16.9^\circ$
      или $196.9^\circ$ (180$^\circ$-вырождение).
\item $\Delta m^2_{21}/\Delta m^2_{31}$: диапазон
      $10^{-5}\ldots 10^{-1}$ при одинаковом loss.
      PDG: $0.0307$.
\item $m_{\beta\beta}$: $0.002$--$0.014$~эВ.
\end{itemize}

\begin{remark}
Нейтринный сектор имеет 12 непрерывных фазовых
параметров, тогда как число наблюдаемых --- 5.
Система недоопределена: любому значению $\Delta m^2$
можно сопоставить точку в 7-мерном многообразии
минимумов без изменения углов. Это означает, что
нейтринные массы в текущей архитектуре не могут
считаться предсказаниями модели.
\end{remark}

% ============================================================
\section{Проверенные гипотезы}
% ============================================================

\begin{table}[H]
\centering
\caption{Одиннадцать независимых тестов модели.}
\label{tab:tests}
\begin{tabular}{clcc}
\toprule
\# & Тест & Результат & Статус \\
\midrule
1 & Космология (арифметика)          & согласована & $\checkmark$ \\
2 & Массы (FN)                       & $10^{-19}\ldots10^{-21}$ & $\checkmark$ \\
3 & CKM                              & $0.1\ldots 5.2\%$ & $\checkmark$ \\
4 & PMNS                             & $0.3\ldots 1.5\%$ & $\checkmark$ \\
5 & Ранг якобиана                    & 4 + 4 & $\checkmark$ \\
6 & Разрыв Гессиана                  & 4 + 4 & $\checkmark$ \\
7 & Калибровочные направления        & $\ge 9$ из 58 & $\checkmark$ \\
8 & Устойчивость нейтринных предсказаний & только NO & частично \\
9 & Альтернативные структуры $M_R$   & не улучшают & $\times$ \\
10 & Унификация $\sigma_R$ с penalty & подгонка, не предсказание & $\times$ \\
11 & Дискретные $\mathbb{Z}_3$-фазы  & 0 из 531\,441 с loss $< 10^{-4}$ & $\times$ \\
\bottomrule
\end{tabular}
\end{table}

% ============================================================
\section{Сравнение с аналогичными моделями}
% ============================================================

MEML 75 принадлежит к классу моделей, в которых
иерархия масс возникает из геометрии дополнительных
измерений.

\begin{itemize}[leftmargin=1.5em,itemsep=0.15em]
\item \textbf{Магнетизированные орбифолды}
      (Hoshiya, Kikuchi, Kobayashi, Uchida):
      киральные фермионы как нулевые моды
      в магнитном поле.
\item \textbf{Модулярные модели} (Feruglio,
      Kobayashi, Ratz и др.): юкавские связи
      как модулярные формы.
\item \textbf{Орбифолдные GUT} (Haba, Shimizu,
      Nomura): нарушение большой группы
      граничными условиями.
\end{itemize}

По соотношению ``наблюдаемые / физические
параметры'': Panin~(2026) --- 9 на 4 (отношение 0.44);
модулярные модели --- 0.6--1.0;
MEML 75 --- 14 на 14 (отношение 1.00). По экономичности
MEML 75 уступает лидерам, но выигрывает в полноте
охвата: одновременно описывает массы, CKM, PMNS
и нейтринные наблюдаемые.

% ============================================================
\section{Итоговый статус модели}
% ============================================================

\subsection{Что модель делает}

\begin{itemize}[leftmargin=1.5em,itemsep=0.15em]
\item Воспроизводит 6 отношений масс фермионов
      с точностью $10^{-19}\ldots10^{-21}$.
\item Воспроизводит CKM с точностью $0.1\ldots 5.2\%$.
\item Воспроизводит PMNS с точностью $0.3\ldots 1.5\%$,
      включая CP-нарушение $|J|$ с точностью $0.9\%$.
\item Устойчиво предсказывает иерархию масс
      нейтрино (NO).
\item Даёт согласованный космологический блок
      (аксионные параметры).
\end{itemize}

\subsection{Что модель не делает}

\begin{itemize}[leftmargin=1.5em,itemsep=0.15em]
\item Не выводит калибровочный сектор СМ.
\item Не выводит хиггсовский сектор.
\item Не квантует гравитацию.
\item Не объясняет инфляцию и бариогенезис.
\item Не является теорией всего.
\item Не предсказывает устойчиво
      $\Delta m^2_{21}/\Delta m^2_{31}$,
      $\delta_{CP}$, $\Sigma m_\nu$, $m_{\beta\beta}$.
\item Не допускает дискретизации фаз
      к сетке $\mathbb{Z}_3$.
\end{itemize}

% ============================================================
\section{Дорожная карта}
% ============================================================

\subsection{Ближний горизонт (1--3 месяца)}

\paragraph{1.1. Проверка 6D-аномалий.}
Для полноценного шестимерного поля необходимо
проверить сокращение калибровочных, гравитационных
и смешанных аномалий на $T^2/\mathbb{Z}_3$. Это
единственный крупный невыполненный тест. Если
аномалии не сокращаются, модель требует расширения
спектра.

\paragraph{1.2. Дополнительные структуры $M_R$.}
Проверить несимметричную $M_R$ с независимыми
диагональными и недиагональными модулями.
Цель --- либо улучшить согласие с
$\Delta m^2_{21}/\Delta m^2_{31}$, либо окончательно
показать предел FN-анзаца.

\paragraph{1.3. Оформление результатов.}
Публикация настоящей работы на arXiv как
феноменологической модели.

\subsection{Средний горизонт (6--12 месяцев)}

\paragraph{2.1. Калибровочный сектор.}
Попытка вывести $SU(3)\times SU(2)\times U(1)$
из калибровочной группы более высокого ранга
($SU(5)$, $SO(10)$, $E_6$) на орбифолде.
Основная сложность --- воспроизвести точные
значения констант связи при электрослабой шкале.

\paragraph{2.2. Хиггсовский сектор.}
Включить хиггсовское поле как динамическую
переменную. Цель: $m_H \approx 125$~ГэВ,
$v = 246$~ГэВ, $\lambda \approx 0.13$.

\paragraph{2.3. Лептогенезис.}
Численное решение уравнений Больцмана для
эволюции барионной асимметрии.

\subsection{Дальний горизонт (2--5 лет)}

\paragraph{3.1. Квантовые поправки.}
Вычисление однопетлевых диаграмм для масс
и углов смешивания. Проверка устойчивости
древесных результатов.

\paragraph{3.2. Гравитационный сектор.}
Включение гравитации как динамической степени
свободы. Это превращает модель в полноценную
теорию гравитации с локализованной материей.

\paragraph{3.3. Экспериментальные предсказания.}
Проверка аксионных параметров ($m_a = 28.22$~мкЭВ,
$\nu_a = 6.824$~ГГц) в ADMX, CASPEr. Поиск
отклонений от СМ на LHC, FCC, DUNE, Hyper-K.

% ============================================================
\section{Заключение}
% ============================================================

Мы представили финальную ревизию модели MEML 75 ---
феноменологической параметризации фермионного
сектора СМ на орбифолде $T^2/\mathbb{Z}_3$.

\textbf{Основные результаты:}
\begin{itemize}[leftmargin=1.5em,itemsep=0.15em]
\item 6 отношений масс: точность $10^{-19}\ldots10^{-21}$.
\item CKM: $0.1\ldots 5.2\%$.
\item PMNS: $0.3\ldots 1.5\%$, включая CP.
\item Число физических параметров в фермионном
      секторе: 14 (6 + 4 + 4).
\item Устойчивое предсказание: иерархия масс
      нейтрино (NO).
\item Космология: $f_a = 1.99 \times 10^{11}$~ГэВ,
      $m_a = 28.22$~мкЭВ, $\nu_a = 6.824$~ГГц.
\end{itemize}

\textbf{Одиннадцать независимых тестов} показали,
что модель \emph{воспроизводит} фермионный сектор
СМ как \emph{феноменологическая параметризация},
но не имеет предсказательной силы за пределами
иерархии масс нейтрино.

\textbf{Открытые задачи} (дорожная карта):
6D-аномалии, калибровочный и хиггсовский секторы,
лептогенезис, квантовые поправки, гравитация.

Модель следует рассматривать как
\emph{феноменологический инструмент} для
систематизации фермионных наблюдаемых,
а не как фундаментальную теорию.

\vspace{0.5cm}

\section*{Благодарности}

Автор благодарит участников обсуждений за
критические замечания, способствовавшие
уточнению формулировок и устранению
ошибочных заявлений в ранних ревизиях.

\vspace{0.3cm}

\begin{thebibliography}{99}

\bibitem{ArkaniHamed1999}
N.~Arkani-Hamed, M.~Schmaltz,
\textit{Hierarchies without symmetries from extra dimensions},
Phys. Rev. D \textbf{61}, 033005 (2000), arXiv:hep-ph/9903417.

\bibitem{Planck2018}
Planck Collaboration,
\textit{Planck 2018 results. VI.},
Astron. Astrophys. \textbf{641}, A6 (2020), arXiv:1807.06209.

\bibitem{PecceiQuinn1977}
R.~D. Peccei, H.~R. Quinn,
\textit{CP Conservation in the Presence of Pseudoparticles},
Phys. Rev. Lett. \textbf{38}, 1440 (1977).

\bibitem{PDG2024}
R.~L. Workman et al. (PDG),
\textit{Review of Particle Physics},
Prog. Theor. Exp. Phys. \textbf{2022}, 083C01 (2022),
2024 update.

\bibitem{KamLANDZen2024}
KamLAND-Zen Collaboration,
\textit{Search for Majorana Neutrinos with the Complete
KamLAND-Zen Dataset}, arXiv:2406.11438 (2024).

\bibitem{ADMX2025}
ADMX Collaboration,
Phys. Rev. Lett. \textbf{134}, 021002 (2025).

\bibitem{Feruglio2017}
F.~Feruglio,
\textit{Are neutrino masses modular flavour models?},
Eur. Phys. J. C \textbf{77}, 713 (2017), arXiv:1706.08749.

\bibitem{Kobayashi2019}
T.~Kobayashi et al.,
\textit{Modular $S_3$-invariant flavor model in SU(5) GUT},
Prog. Theor. Exp. Phys. \textbf{2020}, 053B05 (2020),
arXiv:1906.10341.

\bibitem{Panin2026}
A.~A. Panin,
\textit{Geometric origin of fermion mass hierarchy},
(2026), препринт.

\bibitem{Haba2003}
N.~Haba, Y.~Shimizu, T.~Yamashita,
\textit{Quark and lepton masses from 5D SO(10)},
JHEP \textbf{0306}, 006 (2003), arXiv:hep-ph/0304211.

\end{thebibliography}

\appendix

\section{Численные детали}

\subsection{Алгоритм}

Фит --- L-BFGS-B, численный градиент
($\varepsilon = 10^{-8}$). Рестарты: 16 (массы),
32 (CKM), 128 (PMNS). Параллелизация через
\texttt{ProcessPoolExecutor}, ограничение
BLAS-потоков до 1.

\subsection{Штраф за $M_R$}

Линейный штраф при $\lambda_{\min} < 10^6$~ГэВ:
\begin{equation}
\mathcal{P}(M_R) = 10^3 + 10^6 \cdot
\frac{10^6 - \lambda_{\min}}{10^6}.
\end{equation}

\subsection{Дискретные фазы}

Полный перебор $3^{12} = 531\,441$ комбинаций
через \texttt{einsum}. Время: $< 1$~с на вариант
при векторизации.

\subsection{Скрипты}

Скрипты расчётов доступны в архиве:
\texttt{meml\_all\_checks.py},
\texttt{meml\_v76\_8\_stability.py},
\texttt{meml\_v77\_0\_mixed.py},
\texttt{meml\_v77\_3\_unified.py},
\texttt{meml\_v77\_5\_discrete\_z3.py}.
Воспроизводимость: \texttt{numpy}~$\ge 1.24$,
\texttt{scipy}~$\ge 1.10$.

\end{document}